\section{Completion of the full-dimension induction}
\label{sec:completion}

Theorem~\ref{thm:terminal-big} treats varieties with terminal
singularities and linearly trivial canonical divisor.  We now pass to
klt Calabi--Yau pairs and complete the full-dimension assertion.  The
boundary perturbation below follows the reduction in
\cite[Section~7]{LP}: an effective representative of a small boundary
perturbation plus the nef divisor allows ordinary log abundance to be
applied to a multiple of that nef divisor.

\begin{proposition}\label{prop:full-all}
Let $n\geq 2$, and assume $(P_j)$ for every $j<n$, together with
Theorems~\ref{input:abundance}
and~\ref{input:models}.  Let $(X,B)$ be a projective klt $\Q$-pair over
$\C$ of dimension $n$, and let $M$ be a nef $\Q$-Cartier divisor on $X$.
Suppose that $T=K_X+B$ is pseudo-effective and that, for some rational
$a>0$,
\[
  (T+aM)\cdot C>0
\]
for every integral curve $C$ through a very general point of $X$.
Then $T+cM$ is big for every rational $c>0$.
\end{proposition}

\begin{proof}
Proposition~\ref{prop:full-positive} proves the assertion when
$\kappa(T)>0$.  Proposition~\ref{prop:reduction-cy} reduces the remaining
case to the following claim in dimension $n$:
\begin{equation}\label{eq:completion-cy-claim}
 \begin{gathered}
 (X,B)\text{ projective klt},\qquad K_X+B\equiv 0,\\
 M\text{ nef and $\Q$-Cartier of full nef dimension}
 \quad\Longrightarrow\quad M\text{ big}.
 \end{gathered}
\end{equation}
We prove this claim first with $B=0$ and then with arbitrary boundary.
We may take a small projective $\Q$-factorialization throughout: the
adjoint and $M$ pull back, full nef dimension is preserved on the
isomorphism locus, and bigness of the pullback implies bigness downstairs.

\medskip
\noindent\emph{The case of zero boundary.}
Suppose that $X$ is klt and $K_X\equiv0$.  Theorem~\ref{input:abundance}
makes $K_X$ semiample, hence $K_X\sim_{\Q}0$.  Indeed, a semiample
numerically trivial divisor defines a morphism with zero-dimensional
image, so a sufficiently divisible multiple is linearly trivial.

Take the global index-one cover
\[
  \pi\colon X^{\mathrm{ind}}\longrightarrow X
\]
associated with a trivialization of the smallest positive multiple of
$K_X$ that is linearly trivial.  This is a finite quasi-\'etale cover,
and
\[
  K_{X^{\mathrm{ind}}}=\pi^*K_X\sim0.
\]
The cover is klt.  Since its canonical divisor is Cartier, its log
discrepancies are positive integers, hence at least one; it therefore
has canonical singularities.
These are the standard index-one cover properties for klt varieties;
see~\cite[Definition~5.19 and Proposition~5.20]{KM}.
By the crepant terminalization theorem
\cite[Corollary~1.4.3]{BCHM}, there is a projective birational morphism
\[
  h\colon \widetilde X\longrightarrow X^{\mathrm{ind}}
\]
with $\widetilde X$ terminal and $\Q$-factorial and
$K_{\widetilde X}=h^*K_{X^{\mathrm{ind}}}$.
In particular, $K_{\widetilde X}$ is Cartier and linearly trivial, as
required by Theorem~\ref{thm:terminal-big}.

Set $\widetilde M=(\pi\circ h)^*M$.  This divisor is nef and has full
nef dimension.  To verify the latter assertion, choose a very general
point upstairs outside the exceptional locus of $h$ and above the
very general locus used for $M$.  Every curve through this point has
a curve as its image in $X$, and the projection formula gives strictly
positive $\widetilde M$-degree.  A positive Cartier multiple of
$\widetilde M$ is therefore big by Theorem~\ref{thm:terminal-big}.
Bigness is detected by a dominant generically finite pullback, so $M$
is big on $X$.

\medskip
\noindent\emph{Introducing the boundary.}
Now let $(X,B)$ satisfy the hypotheses of
\eqref{eq:completion-cy-claim}, with $B\ne0$.
Theorem~\ref{input:abundance} gives
\[
  K_X+B\sim_{\Q}0.
\]
Choose a rational $\alpha>0$ sufficiently small that
$(X,(1+\alpha)B)$ remains klt, and put
\[
  T_*=K_X+(1+\alpha)B\sim_{\Q}\alpha B.
\]
Thus $T_*$ is pseudo-effective and $\kappa(T_*)\geq0$.  The full
condition holds for $T_*$ and $M$: for a curve $C$ through a point
outside $\Supp B$, we have $T_*\cdot C=\alpha B\cdot C\geq0$,
whereas $M\cdot C>0$ if that point is very general.
We will show that $T_*+M$ is big without using the conclusion
of the present proposition.

If $\kappa(T_*)>0$, this follows from
Proposition~\ref{prop:full-positive}.  Suppose instead that
$\kappa(T_*)=0$, and apply
Lemma~\ref{lemma:two-nef-model} to $(X,(1+\alpha)B)$ and $M$.
Write $Y$ for its final model, $B_Y$ and $N$ for the transforms of
$B$ and $M$, and
\[
  U=K_Y+(1+\alpha)B_Y.
\]
The output has $U\sim_{\Q}0$ and, for some rational $s>0$, the divisor
$U+sN$ is nef of full nef dimension.  The maps in the construction
extract no divisors, so pushing forward the rational linear equivalence
$K_X+B\sim_{\Q}0$ also gives
\[
  K_Y+B_Y\sim_{\Q}0.
\]
Subtracting these two equivalences shows that
$\alpha B_Y\sim_{\Q}0$.  Since $B_Y$ is effective on a projective
variety, intersection with a sufficiently ample complete-intersection
curve forces $B_Y=0$.  The ordinary pair on $Y$ is therefore klt with
zero boundary and $K_Y\sim_{\Q}0$.
Moreover, $N$ is nef of full nef dimension, because
$U+sN\equiv sN$.  The zero-boundary case just proved makes $N$, and
hence $U+sN$, big.

On a common smooth resolution, with maps $p$ to $X$ and $q$ to $Y$,
the effective exceptional comparison from the reduction is
\[
  p^*(T_*+sM)=q^*(U+sN)+E,
  \qquad E\geq0.
\]
It follows that $T_*+sM$ is big.  Choose a positive rational $\lambda$
with $\lambda<\min\{1,1/s\}$.  Then
\[
  T_*+M
   =\lambda(T_*+sM)+(1-\lambda)T_*+(1-\lambda s)M
\]
is big, since the first summand is big and the other two are
pseudo-effective.  We have thus established bigness of $T_*+M$ in
both cases.

Choose an effective $\Q$-divisor $G$ with
\[
  G\sim_{\Q}T_*+M\sim_{\Q}\alpha B+M.
\]
For a sufficiently small rational $u>0$, the divisor
\[
  \Delta_u=(1-u\alpha)B+uG
\]
is effective and $(X,\Delta_u)$ is klt.  To see this explicitly, take
$u\alpha<1$ and choose $u$ small enough that $(X,B+uG)$ is klt on a
fixed log resolution; then $\Delta_u\leq B+uG$.
Its adjoint satisfies
\[
  K_X+\Delta_u
   \sim_{\Q} K_X+B+u(G-\alpha B)
   \sim_{\Q}uM.
\]
This adjoint is nef, so Theorem~\ref{input:abundance} makes it
semiample.  Consequently $M$ itself is semiample.
The morphism given by a free multiple of $M$ has zero-dimensional
general fiber: a positive-dimensional fiber through a very general
point would contain a curve of $M$-degree zero, contrary to full nef
dimension.  That morphism is therefore generically finite onto its
image, and $M$ is big.  This proves
\eqref{eq:completion-cy-claim} for all klt Calabi--Yau pairs.

Proposition~\ref{prop:reduction-cy} now supplies the case
$\kappa(T)=0$, while Proposition~\ref{prop:full-positive} supplies
$\kappa(T)>0$.  These exhaust the possibilities furnished by the
minimal-model and abundance inputs for a pseudo-effective klt
adjoint, and establish the proposition.
\end{proof}
